% ============================================================
% Document: cws_exposure_backcasting.tex
% Purpose:  Research question, data, and methods for measuring the
%           population served by community water systems downstream of
%           coal mining, 1990-present, including a backcasting procedure
%           for pre-2005 population and an environmental-justice analysis.
% Audience: author, advisor, and economics researchers unfamiliar with
%           the community-water-system context or backcasting methods.
% Author:   EK   Date: 2026-06-17
% ============================================================
\documentclass[11pt]{article}

\usepackage[margin=1in]{geometry}
\usepackage{amsmath,amssymb}
\usepackage{booktabs}
\usepackage{enumitem}

\newcommand{\PWSID}{\texttt{PWSID}}
\newcommand{\popserved}{\ensuremath{S}}
\newcommand{\geopop}{\ensuremath{G}}
\newcommand{\ratio}{\ensuremath{r}}

\title{Measuring the Population Served by Community Water Systems\\
Downstream of Coal Mining, 1990--Present:\\
\large Research Question, Data, and Methods}
\author{Emil Kee-Tui}
\date{\today}

\begin{document}
\maketitle

\begin{abstract}
\noindent
This note sets out a measurement exercise: count the number of people whose
drinking water comes from a community water system (CWS) located downstream of an
active coal mine, track how that count evolved from 1990 to the present, and
characterize who those people are (an environmental-justice, or EJ, analysis). The
central obstacle is that system-level population data do not exist in usable form
before the mid-2000s. We therefore \emph{reconstruct} pre-2005 population using a
transparent ratio-scaling procedure that combines (i) the boundaries of each
system's service area, (ii) decennial Census population counts inside those
boundaries, and (iii) the years when a system's reported population \emph{is}
observed, used to calibrate the reconstruction. The document is written for readers
who know applied microeconomics but not the institutional details of U.S. drinking
water or the ``backcasting'' literature; both are introduced from scratch.
\end{abstract}

\section{Research questions}
\label{sec:questions}

Let a \emph{community water system} (CWS) be ``downstream of coal'' in year $t$ if at
least one active coal mine lies upstream of the system's water intake in the
hydrological network (Section~\ref{sec:background} makes this precise). Write
$\mathcal{D}_t$ for the set of such systems in year $t$, and let $\popserved_{i,t}$
be the number of people served by system $i$ in year $t$. The exercise answers four
questions.

\begin{description}[leftmargin=2em,style=nextline]
  \item[Q1. Pre-period change (1990--2005).] How did total exposed population
    $\sum_{i\in\mathcal{D}_t}\popserved_{i,t}$ change between 1990 and 2005? This is
    the period for which $\popserved_{i,t}$ must be reconstructed.
  \item[Q2. Recent change (2005--present).] How did the same total change between
    2005 and the present, where system-level population is largely observed?
  \item[Q3. No-coal counterfactual.] If all coal mining were removed from the
    network, how many people would \emph{no longer} be downstream of coal? Formally,
    recompute $\mathcal{D}_t$ with the mine layer emptied and report the headcount
    relieved. This is a descriptive accounting counterfactual, not a causal
    water-quality claim (Section~\ref{sec:counterfactual}).
  \item[Q4. Environmental justice.] Among the exposed population, what is the
    demographic composition (race/ethnicity, income, age), and how does it compare to
    the non-exposed population? Are disadvantaged groups over-represented downstream of
    coal, and has that changed over time?
\end{description}

The output of the exercise is a set of \emph{counts and distributions}, not a
regression coefficient. This matters methodologically: an aggregate headcount is far
more forgiving of system-level measurement error than a regression control would be,
because idiosyncratic errors across systems tend to cancel in the sum.

\section{Institutional background}
\label{sec:background}

\paragraph{Community water systems and \PWSID.}
A community water system is a utility that supplies drinking water to the same
resident population year-round (a municipal utility, a rural water district, a
mobile-home park, etc.). Every public water system in the United States carries a
unique identifier, the \PWSID, assigned in the Safe Drinking Water Information System
(SDWIS), the federal regulatory database maintained by the Environmental Protection
Agency (EPA). A key administrative field, \texttt{POPULATION\_SERVED\_COUNT}, records
the estimated resident population the system serves. Two features of this field drive
the entire methodology:
\begin{enumerate}[leftmargin=2em]
  \item It is a \emph{snapshot}, not a time series. EPA overwrites the value as
    systems file updates, so a single download yields one number per system, dated
    to the download, with no history.
  \item It is a \emph{reported} quantity (the population the utility claims to serve),
    which need not equal the residential population living inside the system's
    service-area boundary --- some residents use private wells, and some systems serve
    transient or commercial demand. The gap between the two is central below.
\end{enumerate}

\paragraph{Watersheds and the ``downstream-of-mine'' classification.}
Surface-water flow is organized by the U.S. Geological Survey into nested watershed
units; the finest standard unit used here is the twelve-digit hydrologic unit code
(HUC12), a sub-watershed averaging roughly $100$\,km$^2$. Each HUC12 records the
identity of the single HUC12 it drains into (its \texttt{tohuc}), so the HUC12 layer
forms a directed flow network. A CWS is matched to the HUC12 containing its intake.
Using mine locations (from the Mine Safety and Health Administration, MSHA) placed on
the same network, a HUC12 --- and the systems drawing from it --- is classified as
\emph{downstream of a mine} in year $t$ if an active mine sits in any HUC12 that flows
into it. \textbf{The author already possesses the resulting year-by-year roster
$\mathcal{D}_t$ for 1985--2025}, including which systems are downstream of which mines.
This roster is an \emph{input} to the present exercise, not something it constructs.

\paragraph{Why this is hard.}
The treatment roster $\mathcal{D}_t$ is known annually, but the quantity we wish to
sum over it --- $\popserved_{i,t}$ --- is observed only at a handful of recent dates
(Section~\ref{sec:data}). The methodological core (Section~\ref{sec:methods}) is a
procedure to fill $\popserved_{i,t}$ backward to 1990.

\section{Data}
\label{sec:data}

\subsection{Inputs the author already holds}
\begin{itemize}[leftmargin=2em]
  \item \textbf{Downstream roster $\mathcal{D}_t$, 1985--2025}, annual, at the
    \PWSID\ level, derived from MSHA mine activity and the HUC12 flow network. Includes
    system existence (entry/exit) by year.
  \item \textbf{Service-area polygons} for CWS: the geographic boundary each system
    nominally serves. These are the linchpin that lets us move from coarse county
    proxies to system-specific population (Section~\ref{sec:apportion}).
  \item One modern \texttt{SDWA\_PUB\_WATER\_SYSTEMS} snapshot (current SDWIS
    inventory) as a present-day population anchor.
\end{itemize}

\subsection{Census geography and demography}
The residential population and demographics inside a service polygon come from the
Census Bureau:
\begin{itemize}[leftmargin=2em]
  \item \textbf{Tract/block boundary polygons and 100\% population counts} at the
    decennial censuses \textbf{1990, 2000, 2010, 2020}. Nationwide digital boundary
    files begin with the 1990 TIGER/Line release; \emph{this is why the study window
    opens in 1990 and not earlier}. (1980 tract polygons exist only for already-tracted
    metropolitan areas via IPUMS NHGIS and are excluded.)
  \item \textbf{American Community Survey (ACS)} five-year tract estimates of
    population and demographics, available for every year from $\approx$2009 onward,
    filling the intercensal years in the recent period.
  \item \textbf{Population Estimates Program (PEP)} annual \emph{county} population,
    available continuously, used to give the system-level series an annual cadence
    between decennial points (Section~\ref{sec:cadence}).
\end{itemize}

\subsection{Reported population-served anchors}
``Anchors'' are dates at which we directly observe a system's reported population,
keyed to \PWSID\ so they attach to the roster. Table~\ref{tab:anchors} inventories
the usable sources; the key distinction is whether a source is linkable to \PWSID\
(usable as a per-system anchor) or only a de-identified sample (usable for
distributional validation, not anchoring).

\begin{table}[ht]
\centering
\caption{Sources of reported CWS population served, by \PWSID-linkability.}
\label{tab:anchors}
\small
\begin{tabular}{@{}p{4.5cm}cp{2.1cm}p{5.4cm}@{}}
\toprule
Source & Year(s) & \PWSID-linked? & Content / role \\
\midrule
SYR2 occurrence inventory & 1998--2005 & Yes & Embedded inventory, population served. Carries the early-window anchor. \\
SYR3 occurrence inventory & 2006--2011 & Yes & Per-chemical files (e.g.\ \texttt{barium.txt}) carry \PWSID, \texttt{Retail Population Served}, and \texttt{Adjusted Total Population Served}, dated by sample collection date --- a multi-year anchor bridging SYR2 and the 2010/2011 freezes. \\
CWSS 2006 \texttt{cwssframe} & 2006 & Yes & 51{,}081 CWS; raw \texttt{retailpop}/\texttt{totalpop}. \\
SDWIS 2010 freeze & 2010 & Yes & 93{,}526 CWS; raw \texttt{RetPopSrvd}; also existence dates, facility counts, wholesaler flag, county. \\
SDWIS 2011 freeze & 2011 & Yes & raw \texttt{RetPopSrvd}. \\
\texttt{SDWA\_PUB\_WATER\_SYSTEMS} & current & Yes & present-day snapshot. \\
Buyers/Sellers crosswalk & 2010 & Yes & 39{,}885 buyer/seller \PWSID\ pairs for wholesale de-duplication. \\
\midrule
CWSS 1995 / CWSS 2000 & 1995/2000 & \textbf{No} & De-identified stratified samples ($\sim$2{,}000 / 1{,}250 systems); validation/distribution only. \\
SDWIS ``2005 freeze'' & 2005 & \textbf{No} & 87-row pivot summary, no system records. \\
\bottomrule
\end{tabular}
\end{table}

\subsection{The binding data gap}
There is \emph{no} \PWSID-linkable system-level population before SYR2 ($\approx$1998).
The 1995 and 2000 Community Water System Surveys (CWSS) are de-identified samples and
cannot be attached to particular systems. Consequently the 1990--1998 portion of Q1 is
the hardest: it must be reconstructed without a direct system-level anchor, relying on
the Census residential population inside each service polygon. This limitation is made
explicit in every result for that sub-window.

\section{Methods}
\label{sec:methods}

\subsection{The core identity}
For any system $i$ with service polygon $A_i$, define two distinct population concepts:
\begin{align}
\geopop_{i,t} &= \text{residential population living inside } A_i \text{ in year } t,
  \quad\text{(from Census)} \\
\popserved_{i,t} &= \text{population the system reports serving in year } t.
  \quad\text{(from SDWIS/CWSS anchors)}
\end{align}
The reconstruction rests on a single, transparent identity linking them through a
system-specific \emph{capture ratio} $\ratio_{i,t}$:
\begin{equation}
\boxed{\;\popserved_{i,t} \;=\; \ratio_{i,t}\,\cdot\,\geopop_{i,t},
\qquad \ratio_{i,t} \equiv \popserved_{i,t}/\geopop_{i,t}.\;}
\label{eq:identity}
\end{equation}
$\geopop_{i,t}$ is observable in every Census year directly from the polygon and the
tract data; $\ratio_{i,t}$ is observable only in anchor years. The strategy is:
estimate $\geopop_{i,t}$ everywhere, estimate $\ratio_{i,t}$ where anchors exist,
carry $\ratio$ to non-anchor years under an explicit assumption, and multiply.

For readers new to this style of estimate, the logic mirrors how demographers
``backcast'' missing population: a slowly moving, mechanically determined quantity
(here, residents inside a fixed boundary, $\geopop$) serves as a deterministic
backbone, and the comparatively small, system-specific deviation from it (here,
$\ratio$) is the only thing that must be modeled. Recent work casts exactly this kind
of intercensal population gap-filling as a prediction problem and shows the backbone
carries almost all of the explanatory power, with machine learning adding only a
modest refinement on top (Section~\ref{sec:ml}).

\subsection{Step 1 --- Spatial apportionment of $\geopop$}
\label{sec:apportion}
Because a Census tract or block rarely coincides with a service polygon, we compute
$\geopop_{i,t}$ by \emph{population-weighted areal interpolation}: intersect polygon
$A_i$ with the Census blocks of decennial year $t$, and sum each block's population
weighted by the share of that block's population falling inside $A_i$:
\begin{equation}
\geopop_{i,t} \;=\; \sum_{b\,\in\,\text{blocks}} \texttt{pop}_{b,t}\,\cdot\,
  w_{b,i}, \qquad
  w_{b,i} = \frac{\text{population of } b \text{ inside } A_i}{\text{population of } b}.
\label{eq:apportion}
\end{equation}
Blocks (the finest Census unit, with 100\% counts at each decennial) make $w_{b,i}$
accurate; where only tracts are available (e.g.\ ACS years) the same formula is
applied at tract resolution. Crucially, because the \emph{polygon $A_i$ is fixed},
the recurring problem that Census tract boundaries change every decade simply
disappears: we re-intersect each decade's geography with the same target polygon, and
never need a tract-to-tract crosswalk. Demographic composition (race/ethnicity,
income, age shares) is apportioned into $A_i$ by the identical weights $w_{b,i}$,
which feeds the EJ analysis (Section~\ref{sec:ej}).

\subsection{Step 2 --- Annual cadence for $\geopop$}
\label{sec:cadence}
Equation~\eqref{eq:apportion} yields $\geopop_{i,t}$ only at Census years
($1990,2000,2010,2020$; plus ACS from $\approx$2009). To obtain a value for every
year $t$, we grow $\geopop_{i,t}$ forward year by year from the most recent decennial
polygon apportionment, applying the county's \emph{own annual growth rate} from the
Population Estimates Program (PEP). Concretely, whatever the county population grows by
between $t-1$ and $t$, the service-area population is assumed to grow by the same
percentage: if the county grew $1\%$ from 1990 to 1991, the system's $\geopop$ is
raised $1\%$ over the same year. Growth is always taken from $t-1$ to $t$, chained
forward from the opening decennial anchor.

\paragraph{Notation.}
Fix a system $i$ and let $c(i)$ be the county whose annual PEP series we borrow; in
practice $c(i)$ is the county containing the centroid of $A_i$, or, when $A_i$
straddles counties, the population-weighted blend
$P_{c(i),t}=\sum_{c}\theta_{i,c}\,P_{c,t}$ with weights $\theta_{i,c}$ equal to the
share of $\geopop_{i,t_0}$ drawn from county $c$ at the opening decennial. Write
$P_{c(i),t}$ for that (possibly blended) county population in year $t$, observed
annually. Let $t_0$ be the most recent decennial Census year at or before $t$
(e.g.\ $t_0=1990$ for $1990\le t<2000$), at which the polygon apportionment of
Equation~\eqref{eq:apportion} gives the anchor $\geopop_{i,t_0}$.

\paragraph{Chained county-growth rule.}
The annual value is defined recursively by applying the county's gross growth factor
$P_{c(i),t}/P_{c(i),t-1}$ to the previous year's estimate, starting from the decennial
anchor:
\begin{equation}
\boxed{\;
\hat\geopop_{i,t}\;=\;\hat\geopop_{i,t-1}\,\cdot\,
   \frac{P_{c(i),t}}{P_{c(i),t-1}},
\qquad \hat\geopop_{i,t_0}\;=\;\geopop_{i,t_0}.
\;}
\label{eq:cadence}
\end{equation}
The recursion telescopes to a closed form: the cumulative product of annual growth
factors collapses to the ratio of county populations at the two endpoints, so
\begin{equation}
\hat\geopop_{i,t}\;=\;\geopop_{i,t_0}\,\cdot\,
   \frac{P_{c(i),t}}{P_{c(i),t_0}}
\;=\;\kappa_{i,t_0}\,P_{c(i),t},
\qquad
\kappa_{i,t_0}\equiv\frac{\geopop_{i,t_0}}{P_{c(i),t_0}}.
\label{eq:cadence-closed}
\end{equation}
That is, the rule holds the system's \emph{coverage ratio} --- its share
$\kappa_{i,t_0}$ of the county population --- fixed at the opening decennial value, and
lets the service-area population ride the county's annual trajectory thereafter. The
single substantive assumption is that $\kappa$ is constant within the decade: the
polygon captures the same fraction of its county every year. This is exactly the
degenerate case of a full two-endpoint benchmarking rule in which the closing-decennial
coverage ratio happens to equal the opening one; here we impose it rather than
estimate it.

\paragraph{Re-anchoring and boundary conventions.}
At each decennial we reset the level to the freshly observed apportionment
($\hat\geopop_{i,t_1}\!\leftarrow\!\geopop_{i,t_1}$) and resume chaining from there;
because the forward path is not benchmarked to the closing decennial, it will in
general \emph{not} arrive at $\geopop_{i,t_1}$ on its own, so this reset introduces a
discontinuity at the decade boundary equal to the decade's realized drift in $\kappa$.
Two further conventions complete the rule:
\emph{(i)} where an annual ACS apportionment of $\geopop_{i,t}$ is itself available
(post-$\approx$2009), it is used directly and the chained PEP growth is reserved for
the pre-ACS gaps; \emph{(ii)} the same forward rule covers years after the latest
decennial with no special case --- post-2020 simply continues
$\hat\geopop_{i,t}=\kappa_{i,2020}\,P_{c(i),t}$.

\subsection{Step 3 --- The reported layer: calibrating $\ratio$}
At each anchor year $a$ (Table~\ref{tab:anchors}) we observe $\popserved_{i,a}$ and
can compute $\geopop_{i,a}$, hence $\ratio_{i,a} = \popserved_{i,a}/\geopop_{i,a}$.
The ratio absorbs everything that makes reported service differ from polygon residence
(private-well penetration, transient demand, imprecise boundaries). We then form an
estimate $\hat{\ratio}_{i,t}$ for all years:
\begin{itemize}[leftmargin=2em]
  \item \textbf{One anchor:} hold $\ratio$ flat at its anchor value; all time
    variation then comes from $\geopop$.
  \item \textbf{Multiple anchors:} interpolate $\ratio$ linearly between them and hold
    it flat outside the anchored range. (Interpolating the ratio, not the level, lets
    a system's captured share drift with annexation or consolidation while $\geopop$
    carries the demographic trend.)
  \item \textbf{No anchor at all:} impute $\hat{\ratio}$ from the median ratio of
    same-type, same-size systems in the same county; flag these as least reliable.
\end{itemize}
Before any system that buys water from another is counted, the Buyers/Sellers
crosswalk is applied so that population served wholesale is not double-counted across
the selling and buying systems.

\subsection{Step 4 --- Two estimands and the final estimator}
We report exposure under \emph{both} population concepts, which bracket the truth:
\begin{align}
\text{Residential exposure:} &\quad E^{G}_t = \sum_{i\in\mathcal{D}_t}\geopop_{i,t},\\
\text{Reported-served exposure:} &\quad E^{S}_t =
  \sum_{i\in\mathcal{D}_t}\hat{\ratio}_{i,t}\,\geopop_{i,t}.
\end{align}
$E^{G}_t$ counts residents living inside the service areas of downstream systems ---
arguably the right notion of who is \emph{exposed} to the system's water. $E^{S}_t$
matches the utilities' own reported caseload. Reporting both is honest about the
$\ratio$ wedge and gives the reader a range.

\subsection{Operationalizing Q1 and Q2}
\label{sec:q1q2}
\textbf{Q2 (2005--present)} is well-anchored: $\geopop$ comes from the 2010/2020
decennials and ACS, and $\ratio$ from the 2006, 2010, 2011, and current anchors. The
change $E_{\text{present}}-E_{2005}$ is a direct, defensible difference.

\textbf{Q1 (1990--2005)} is partly anchored and partly extrapolated, and we are
explicit about which is which:
\begin{itemize}[leftmargin=2em]
  \item \emph{1998--2005} is anchored by SYR2 (and SYR3/2006 just beyond, continuing
    through 2011), with $\geopop$ from the 1990 and 2000 decennials.
  \item \emph{1990--1997} has \emph{no} system-level anchor. Here exposure is driven
    by $\geopop$ (the 1990 and 2000 decennial apportionments and their interpolation)
    times $\hat{\ratio}$ held flat at its earliest observed (SYR2-era) value. Any
    1990--1997 trend is therefore essentially the Census residential trend inside the
    fixed downstream service areas, and is reported with that caveat stated.
\end{itemize}
Because the downstream roster $\mathcal{D}_t$ is known annually, the change in exposed
population correctly reflects systems entering and leaving the downstream set
(the extensive margin), not only changes in per-system population.

\subsection{Q3 --- The no-coal counterfactual}
\label{sec:counterfactual}
``If all coal disappeared'' is operationalized mechanically. Let $\mathcal{D}_t$ be
the factual downstream set (mines active as observed) and $\mathcal{D}^{0}_t$ the set
recomputed after \emph{emptying the mine layer} from the flow network. With no active
mines, no system is downstream of a mine, so $\mathcal{D}^{0}_t=\varnothing$ and the
relieved population is simply
\begin{equation}
R_t \;=\; \sum_{i\in\mathcal{D}_t}\hat{\popserved}_{i,t} \;-\;
          \sum_{i\in\mathcal{D}^{0}_t}\hat{\popserved}_{i,t}
        \;=\; E_t .
\end{equation}
Two refinements make this more informative than a single number:
\begin{enumerate}[leftmargin=2em]
  \item \textbf{Intensity weighting.} Rather than a binary downstream indicator,
    characterize exposure by the number of upstream mines, upstream production, or
    upstream coal sulfur content, and report the relieved population by intensity tier
    (e.g.\ how many people are downstream of high-sulfur vs.\ low-sulfur coal).
  \item \textbf{Active vs.\ legacy coal.} Removing only \emph{active} mines differs
    from removing all coal influence, since abandoned mine lands can continue to affect
    water. We compute the counterfactual for active mines (primary) and note the
    legacy variant as a bound.
\end{enumerate}
This is descriptive accounting --- a count of who is currently downstream --- and is
explicitly \emph{not} a claim that water quality would improve by a measured amount.

\subsection{Q4 --- Environmental-justice analysis}
\label{sec:ej}
For each downstream system, the same polygon weights $w_{b,i}$ used for population are
applied to Census tract demographic shares (race/ethnicity, median household income,
age, and poverty status) to obtain the demographic composition of the population
served. We then:
\begin{enumerate}[leftmargin=2em]
  \item Compare the exposed population's composition to that of the non-exposed CWS
    population (and to the national population) at each Census year, testing whether
    disadvantaged groups are over-represented downstream of coal.
  \item Track these gaps over 1990--2020 to see whether the burden's demographic
    profile shifted.
  \item Summarize disproportionality with standard EJ measures (population-weighted
    exposure shares by group; optionally a concentration/dissimilarity index).
\end{enumerate}
Because demographic \emph{composition} is far more stable and better measured than the
population \emph{level}, the EJ results are the most robust output of the exercise and
depend least on the backcasting assumptions.

\subsection{Validation}
\label{sec:validation}
Two leave-one-out exercises discipline the procedure, mimicking its actual use:
\begin{itemize}[leftmargin=2em]
  \item \textbf{Leave-one-anchor-out (the $\ratio$ layer).} Withhold a reported anchor
    (e.g.\ 2006), predict it as $\hat{\ratio}\cdot\geopop$ from the remaining anchors,
    and compare to the held-out truth --- at the aggregate (what we report) and
    per-system. Aggregates should be accurate even where individual systems are noisy.
  \item \textbf{Leave-one-decade-out (the $\geopop$ backbone).} Predict the 2000
    polygon apportionment by PEP-scaling from 1990 and 2010, and compare to the true
    2000 block apportionment. This tests the scaling assumption \emph{independently of
    any reported figure} --- something the service polygons newly make possible.
  \item \textbf{Ratio-drift diagnostic.} For multi-anchor systems, report the
    dispersion of $\ratio_{i,a}$ across anchors; its size directly measures how much
    the constant-/interpolated-$\ratio$ assumption strains, rather than assuming it
    away.
\end{itemize}

\subsection{Optional machine-learning robustness}
\label{sec:ml}
The procedure above is deliberately transparent and, for aggregate headcounts, likely
adequate. As a robustness check we can replace the simple $\ratio$ rule and the
PEP-shape interpolation with a supervised-learning predictor (e.g.\ a neural network
or gradient-boosted trees) trained on the anchor years to map exogenous features
(Census county/place population, demographics, system type, source water) to
population served, in the spirit of the demographic backcasting literature. We
deliberately \emph{exclude} SDWA violation and enforcement variables from the feature
set, so that the exposure measure is not mechanically a function of the regulatory
outcomes studied elsewhere in the project. Given that the deterministic backbone
already explains the bulk of the variation, the machine-learning layer is expected to
refine rather than redefine the estimates, and is reported as a sensitivity, not the
headline.

\section{Threats to validity and limitations}
\label{sec:threats}
\begin{enumerate}[leftmargin=2em]
  \item \textbf{Pre-1998 extrapolation.} No system-level anchor exists before SYR2, so
    1990--1997 exposure is a Census-residential extrapolation with $\ratio$ held flat.
    We report it as such and lean on the better-anchored 1998--2005 segment for Q1.
  \item \textbf{1990 polygon floor.} Nationwide tract/block polygons start in 1990;
    the study cannot extend earlier without metro-only, lower-quality 1980 geography.
  \item \textbf{Service-polygon accuracy.} If a polygon poorly delineates the true
    service area, $\geopop$ is mismeasured; the $\ratio$ calibration partly corrects
    this, and the ratio-drift diagnostic flags the worst cases.
  \item \textbf{Capture-ratio non-stationarity.} Holding/interpolating $\ratio$
    assumes a roughly stable captured share; annexation and consolidation violate this.
    Measured drift bounds the error.
  \item \textbf{Reporting comparability.} ``Population served'' is defined slightly
    differently across SYR2, CWSS, and the freezes; definitional drift can masquerade
    as real change and is checked at overlapping years.
  \item \textbf{Aggregation cuts both ways.} Headcounts are robust because errors
    cancel; any \emph{system-level} statement (e.g.\ a specific town) is far less
    reliable than the totals.
\end{enumerate}

\section*{Summary}
\addcontentsline{toc}{section}{Summary}
The known annual roster of systems downstream of coal mining ($\mathcal{D}_t$) is
combined with service-area polygons and Census data to reconstruct, system by system,
the population each serves, via the identity
$\popserved = \ratio\cdot\geopop$: Census apportionment supplies the residential
backbone $\geopop$ for every year back to 1990, and the recent reported anchors
(SYR2, SYR3, CWSS 2006, the 2010/2011 freezes, and a current snapshot) calibrate the
capture ratio $\ratio$. Summing over $\mathcal{D}_t$ yields the exposed-population time series
behind Q1 and Q2; emptying the mine layer yields the no-coal counterfactual of Q3; and
apportioning Census demographics into the same polygons yields the EJ analysis of Q4.
The design's honest weak point is the unanchored 1990--1997 window; its strongest
output is the demographic composition, which barely depends on the backcasting
assumptions at all.

\end{document}
